\section{问题二模型建立及求解}

\subsection{单星多目标巡游模型的建立}

问题二把第 4.3 节的观测单元首尾相接。30 颗目标各对应一条时变可观测管道，R1 要在 30 天内依次穿过尽可能多的管道，每穿过一条就完成一次有效观测，上一次观测结束时的轨道状态就是下一次转移的起点。每条管道大约每天靠近 R1 轨道面一次，于是 30 天内全部目标的观测机会构成一个离散的事件集合。问题二可以概括为：在这些事件中选出一条时间递增、目标互不重复的路径，并为路径上的每一段设计至多两次脉冲，使路径包含的目标最多，累计速度增量最小。路径的选择是组合问题，每段脉冲的设计是连续问题，二者通过中间状态耦合在一起，整体思路如图~\ref{fig:q2abs} 所示。

\begin{figure}[!htb]
  \centering
  \includegraphics[width=0.95\textwidth]{fig9_q2abs.pdf}
  \caption{问题二的抽象：在观测事件之间选取路径并逐段穿过可观测管道}
  \label{fig:q2abs}
\end{figure}

\subsubsection{决策变量定义}

方案由访问序列和各段的机动组成。访问序列记为 $\pi = (\pi_1, \pi_2, \dots, \pi_K)$，$\pi_k$ 是第 $k$ 次观测的目标，$K$ 是观测次数。第 $k$ 次观测之前至多机动两次，这一段的变量与问题一相同，记为
\begin{equation}
  \bm z_k = (t_{k,1},\ t_{k,2},\ \Delta\bm v_{k,1},\ \Delta\bm v_{k,2},\ s_k), \qquad k = 1, 2, \dots, K ,
  \label{eq:q2var}
\end{equation}
其中 $s_k$ 是第 $k$ 次观测的开始时刻。一个完整方案为 $\mathcal P = (\pi;\ \bm z_1, \dots, \bm z_K)$。给定 $\mathcal P$ 后，把全部脉冲按时间排列，用 $\Phi$ 逐段传播，就得到 R1 在 30 天内的整条轨迹。某一段的两次脉冲都取零时，这一段不机动，R1 在上一次观测之后直接飞到下一个目标附近，这种情形对应在同一次飞越中先后观测两颗目标。

\subsubsection{目标函数设置}

按第 4.4 节的两阶段字典序目标，第一阶段使观测到的不同目标数 $K$ 最大，第二阶段在 $K = N^*$ 的方案中最小化累计速度增量
\begin{equation}
  J(\mathcal P) = \sum_{k=1}^{K} \big(\|\Delta\bm v_{k,1}\| + \|\Delta\bm v_{k,2}\|\big) .
  \label{eq:q2obj}
\end{equation}
求解时两个阶段并不分开进行。所有搜索步骤都只接受观测数不减少的候选方案，观测数相同时再比较 $J$，这与式~\eqref{eq:lex1}、\eqref{eq:lex2} 的字典序等价。

\subsubsection{约束条件分析}

约束仍分为时间、脉冲和观测三类，与问题一相比，时间约束和脉冲约束都从单段扩展到了整条时间轴。时间约束要求各段依次进行：第 $k$ 段的第一次机动不早于第 $k-1$ 次观测结束，两次机动间隔不少于 1~h，观测在第二次机动之后开始；相邻两段之间的机动同样要相隔 1~h 以上，最后一次观测在 30 天内结束。脉冲约束要求单次速度增量不超过 30~m/s，累计不超过 300~m/s，在问题二中预算约束可能起作用。观测约束要求第 $k$ 次观测区间内的每一时刻，R1 相对目标 $\pi_k$ 的位置都落在 $\mathcal C_{\pi_k}(t)$ 内，离散化方式与问题一相同。此外，序列中的目标互不重复。

\subsubsection{模型汇总}

记由时间约束和脉冲约束构成的集合为
\begin{equation}
  \mathcal Z_K = \left\{(\bm z_1,\dots,\bm z_K) \;\middle|\;
  \begin{aligned}
    & s_{k-1} + \tau \le t_{k,1},\ \ t_{k,1} - t_{k-1,2} \ge 3600,\ \ t_{k,2} - t_{k,1} \ge 3600, \\
    & t_{k,2} \le s_k,\ \ s_K + \tau \le T,\ \ \|\Delta\bm v_{k,j}\| \le 30~\mathrm{m/s},\ \ J \le 300~\mathrm{m/s}, \\
    & k = 1,\dots,K,\ \ j = 1, 2
  \end{aligned}
  \right\},
  \label{eq:q2set}
\end{equation}
约定 $s_0 + \tau = t_{0,2} + 3600 = 0$，$T = 30$~天。问题二的优化模型为
\begin{equation}
  \operatorname*{lexmax}_{\pi,\ \bm z\in\mathcal Z_K}\ \big(K,\ -J\big) \quad \text{s.t.}\quad
  \begin{aligned}[t]
    & \bm\rho_{\pi_k}(s_k + j;\,\mathcal P) \in \mathcal C_{\pi_k}(s_k + j),\quad j = 0, 1, \dots, 60,\ k = 1, \dots, K, \\
    & \pi_k \ne \pi_l,\quad k \ne l ,
  \end{aligned}
  \label{eq:q2model}
\end{equation}
其中 $\bm\rho_{\pi_k}(t;\mathcal P)$ 是 R1 相对目标 $\pi_k$ 的位置。与问题一的模型~\eqref{eq:q1model} 相比，这里多了离散的序列 $\pi$，而且第 $k$ 次观测的相对位置依赖于前面全部 $k$ 段的机动。

\subsection{优化模型求解}

模型~\eqref{eq:q2model} 的困难来自两方面。离散部分要在约一千个观测事件中选出一条路径，连续部分约有 $9K$ 个变量，而且各段通过中间状态相互影响。我们按“排序、联合、改进”三个层次求解。排序，是在漂移相位模型上用束搜索确定访问顺序和观测日期，再逐段精确求解得到一个初始方案；联合，是固定顺序和时刻，对 30 天计划的全部脉冲做一次联合的序列凸优化；改进，是用大邻域搜索调整访问顺序，每得到一个更好的顺序都重新做联合优化。

\subsubsection{排序：漂移相位模型上的束搜索}

观测事件由第 5.2.1 节的定时方法给出。对每颗目标，求出它穿过 R1 轨道面且 $\beta \le 90^\circ$ 的时刻，每个这样的时刻就是一个观测事件。B7 是例外，它与 R1 的轨道面夹角只有 $0.04^\circ$，到 R1 轨道面的距离始终不超过 50~km，不受穿越条件的限制，在 $\beta \le 90^\circ$ 的时段内每隔 2~h 取一个事件。30 天内共得到 1047 个观测事件。

R1 在事件之间的转移用漂移相位模型估计。把 R1 和目标的位置都投影到 R1 的初始轨道面上，以 R1 的初始平均角速度为参考，定义相位 $\varphi$。要在事件 $(b, s)$ 观测目标 $b$，R1 的相位必须在 $s$ 时刻追上目标的相位。设 R1 当前的漂移率为 $\omega$，相位差为 $\Delta\varphi$，可用时间为 $T_k$ 天，用一次脉冲把漂移率改为 $\omega + \Delta\varphi/T_k$ 即可追上，代价由式~\eqref{eq:drift} 给出，为 $2.84\,|\Delta\varphi/T_k|$~m/s。也可以用两次脉冲，先以较高的漂移率接近，在交会前一天再把漂移率降下来。

交会时允许的漂移率需要单独考虑。切向脉冲使脉冲点成为新轨道的近地点或远地点，此后 R1 在偏离这条拱线 $\delta$ 角的方向上，地心距相对脉冲点变化约 $|\Delta a|(1 - \cos\delta)$。观测发生在目标轨道面与 R1 轨道面的交线方向上，如果交线偏离拱线较远，漂移率越大，R1 到达交线时的高度偏差就越大，最后一次修正的代价也越大。要求这一偏差不超过 $\rho_{\mathrm{tol}}$，结合 $\omega = -\tfrac32 n\,\Delta a/a$，得到交会时允许的最大漂移率
\begin{equation}
  \omega_{\max}(\delta) = \min\left\{\omega_{\mathrm{cap}},\ \frac{3n}{2a}\cdot\frac{\rho_{\mathrm{tol}}}{1 - \cos\delta}\right\} ,
  \label{eq:wmax}
\end{equation}
其中 $\omega_{\mathrm{cap}}$ 是搜索中考虑的最大漂移率，取 $4^\circ$/天；$\rho_{\mathrm{tol}}$ 取 25~km，约为观测距离上限的一半。

拱线方向取 R1 轨道面与赤道面的交线。同步带目标的倾角都很小，它们与 R1 轨道面的交线都靠近 R1 自身的赤道交线，所以把拱线放在这个方向上，可以使尽量多的交会发生在拱线附近。本题中 R1 的倾角为 $4.38^\circ$，30 颗目标里有 18 颗的交线方向与之相差不到 $6^\circ$，连同偏离在 $22^\circ$ 以内的 B2、B15、B28，按式~\eqref{eq:wmax} 都可以用最大漂移率交会；B6、B16、B21、B22、B25、B30 六颗目标的交线偏离超过 $55^\circ$，允许的漂移率不到 $0.75^\circ$/天，只能慢速接近。B7 没有确定的交线方向，按最不利情形 $\delta = 90^\circ$ 处理。

在这一模型上做束搜索\cite{ow1988}。搜索状态为当前时刻、R1 的相位与漂移率、已观测的目标集合和已有序列。每一步从当前状态出发，把未来若干天内所有未观测目标的事件作为候选扩展，按上述方法估计代价，单段代价超过 30~m/s 的扩展直接舍去。扩展后的状态按评价值排序，只保留前 $W$ 个。评价值由三部分组成：
\begin{equation}
  f = J_{\mathrm{est}} + \mu\, t_{\mathrm{day}} + c\,\Lambda_{\mathrm{rem}} ,
  \label{eq:beameval}
\end{equation}
其中 $J_{\mathrm{est}}$ 是已累计的估计代价，$t_{\mathrm{day}}$ 是已用天数，$\Lambda_{\mathrm{rem}}$ 是扫过剩余全部目标还需要的最小相位跨度。第二项避免搜索在一处停留过久，第三项避免搜索只顾一侧、来不及扫过另一侧。漂移相位模型下每次扩展只需代数运算，束宽可以取到数千。

束搜索给出访问顺序和观测日期后，按顺序逐段求解观测单元，每一段就是一个以上一次观测结束状态为初值的问题一，用第 5.2 节的方法求出精确的两脉冲方案，上一段的终端状态接力传给下一段。某一段在排定日期无解时，改用该目标相邻几天的观测事件。逐段求解之后仍缺少的目标，尝试与相邻目标合并为同一次飞越：把两颗目标的观测约束写进同一段，两次观测共用一组机动。

\subsubsection{联合：全计划序列凸优化}

逐段求解只让每一段自己的代价最小，没有考虑这一段的终端状态会给后面几段带来多大负担，所以初始方案不是整体最优。联合优化固定访问顺序、观测时刻和脉冲时刻，把 30 天计划的全部脉冲一起作为变量。

在当前方案的轨迹附近，第 $m$ 次脉冲的修正量 $\delta\Delta\bm v_m$ 对其后任意时刻位置的影响近似是线性的，
\begin{equation}
  \bm r_O(t) \approx \bar{\bm r}_O(t) + \sum_{t_m < t} \bm\Phi_{rv}(t, t_m)\,\delta\Delta\bm v_m ,
  \label{eq:linearize}
\end{equation}
其中 $\bar{\bm r}_O$ 是当前方案的轨迹，$\bm\Phi_{rv}$ 是状态转移矩阵中位置对速度的分块，由精确的二体传播做有限差分得到。第 4.2 节已经指出，可观测集合中两个圆锥和距离上界都是凸约束，只有距离下界不是。我们用半空间 $\hat{\bm g}_k^{\mathrm T}\bm\rho \ge 10$ 代替 $\|\bm\rho\| \ge 10$，$\hat{\bm g}_k$ 取第 $k$ 次观测中相对位置的平均方向。半空间内的点一定满足距离下界，所以这是一个内近似，得到的凸集记为 $\hat{\mathcal C}_b(t)$。60~s 内相对运动近似为直线，在观测区间上均匀取 5 个时刻施加约束，就能控制整段观测。于是联合优化的子问题为
\begin{equation}
  \begin{aligned}
    \min_{\delta\Delta\bm v}\quad & \sum_m \|\Delta\bar{\bm v}_m + \delta\Delta\bm v_m\| \\
    \text{s.t.}\quad & \|\Delta\bar{\bm v}_m + \delta\Delta\bm v_m\| \le 30~\mathrm{m/s},\quad \|\delta\Delta\bm v_m\|_\infty \le \Delta , \\
    & \bm\rho_{\pi_k}(t;\,\delta\Delta\bm v) \in \hat{\mathcal C}_{\pi_k}(t),\quad t \in \{s_k + q\tau/4,\ q = 0,\dots,4\},\ k = 1,\dots,K ,
  \end{aligned}
  \label{eq:socp}
\end{equation}
其中 $\Delta\bar{\bm v}_m$ 是当前方案的脉冲，$\Delta$ 是信赖域半径。由式~\eqref{eq:linearize}，$\bm\rho$ 是 $\delta\Delta\bm v$ 的线性函数，目标函数和全部约束都是二阶锥形式，所以式~\eqref{eq:socp} 是一个二阶锥规划。问题二中它约有 170 个变量，用 CVXPY\cite{cvxpy2016} 调用内点法求解器 Clarabel\cite{clarabel2024} 求解，每次只需几秒。

子问题的解只在线性化成立的范围内可靠，所以采用序列凸优化\cite{mao2016}：用子问题的解更新全部脉冲，再用精确的二体模型重新传播轨迹，在 0.05~s 网格上检验全部观测约束，然后在新的轨迹附近重新线性化，求解下一个子问题。只有通过检验且累计速度增量下降的方案才记为当前最优。信赖域半径每轮乘以 0.9，连续 15 轮没有改进时停止。

\subsubsection{改进：大邻域搜索调整访问顺序}

联合优化不改变访问顺序和观测日期，顺序本身的好坏仍由束搜索的近似代价决定。我们在序列层面做大邻域搜索\cite{shaw1998}。每一步对当前最优方案施加一种邻域动作：交换相邻两次观测的先后；把某次观测改到同一间隙内另一天的事件；把某次观测移到另一个间隙；把连续 2 至 4 次观测整体重排。动作之后，从第一个被改动的位置起逐段重新求解，得到新方案。新方案的观测数不少于当前最优、且累计速度增量更小时才被接受。多个进程并行执行这一过程，共享同一个当前最优方案。搜索结束后，对得到的新顺序再做一次联合优化。

三个层次中，只有排序层用到了漂移相位这一近似模型，联合层和改进层的每一个方案都由精确的二体传播得到，并经过 0.05~s 网格检验，所以最终方案对模型~\eqref{eq:q2model} 严格可行，不依赖于近似模型的精度。

\subsection{结果的分析与验证}

最终方案在 30 天内观测全部 30 颗目标，累计速度增量 176.31~m/s，占预算的 59\%。方案共 58 次脉冲，其中大于 1~mm/s 的有 41 次，单次最大 25.71~m/s，相邻脉冲的最小间隔为 3600~s，最后一次观测在第 29.52 天结束。各次观测的时刻和对应的速度增量见表~\ref{tab:q2plan}。

\begin{table}[!htb]
  \centering
  \caption{问题二观测序列}
  \label{tab:q2plan}
  \zihao{5}\setlength{\tabcolsep}{4pt}
  \begin{tabular}{clcc@{\hspace{14pt}}clcc}
    \toprule
    序号 & 目标 & 观测开始（UTC） & $\Delta v$ / (m/s) & 序号 & 目标 & 观测开始（UTC） & $\Delta v$ / (m/s) \\
    \midrule
1 & B2 & 01-01 14:00:53 & 9.91 & 16 & B12 & 01-20 12:39:42 & 0.10 \\
2 & B3 & 01-02 14:36:43 & 9.21 & 17 & B13 & 01-21 12:35:02 & 2.86 \\
3 & B9 & 01-03 14:28:49 & 0.42 & 18 & B15 & 01-22 12:10:12 & 2.05 \\
4 & B18 & 01-06 12:50:20 & 2.36 & 19 & B14 & 01-23 12:06:41 & 1.52 \\
5 & B30 & 01-08 19:06:54 & 2.60 & 20 & B16 & 01-23 16:50:23 & 0.00 \\
6 & B26 & 01-09 15:03:55 & 8.16 & 21 & B19 & 01-24 12:01:26 & — \\
7 & B17 & 01-12 14:21:42 & 6.01 & 22 & B20 & 01-24 12:12:26 & 9.15 \\
8 & B8 & 01-13 14:07:05 & 1.68 & 23 & B21 & 01-24 17:01:20 & 25.71 \\
9 & B5 & 01-14 13:51:47 & 1.50 & 24 & B23 & 01-25 12:00:09 & 20.74 \\
10 & B1 & 01-15 13:38:06 & 1.16 & 25 & B22 & 01-26 17:48:44 & 2.47 \\
11 & B4 & 01-17 13:06:36 & 0.36 & 26 & B24 & 01-27 11:29:20 & 1.62 \\
12 & B11 & 01-18 12:56:32 & 3.27 & 27 & B25 & 01-27 15:36:18 & 16.21 \\
13 & B6 & 01-18 16:40:13 & 10.71 & 28 & B27 & 01-28 11:34:09 & 18.30 \\
14 & B7 & 01-18 21:45:36 & 10.79 & 29 & B29 & 01-29 13:07:58 & 2.38 \\
15 & B10 & 01-19 12:46:40 & 1.54 & 30 & B28 & 01-30 12:25:14 & 3.52 \\
    \midrule
    \multicolumn{8}{l}{合计 176.31~m/s；B19 与 B20 在同一次飞越中观测，B19 不单独机动。} \\
    \bottomrule
  \end{tabular}
\end{table}

R1 的巡游路线如图~\ref{fig:q2route}(a) 所示。R1 从队列中部出发，先向西漂移，依次观测西侧的 B2、B3、B9、B18，在第 8 天前后到达最西端的 B30，然后折返向东，扫过其余目标，最后一次观测在 R1 初始位置以东约 $18^\circ$ 处。R1 从队列中间出发，两侧都有目标，一次折返是必需的；折返之后路线一直沿相位单向推进，没有再出现往返。

图~\ref{fig:q2route}(b) 给出各次观测的速度增量。多数观测的代价在 3~m/s 以内，中位数为 2.53~m/s，这些目标的交线方向都靠近公共拱线，R1 以较高的漂移率依次交会，每次只需很小的修正。代价较高的几次观测可以分为三类。第一类是 B6、B21、B25 这些交线偏离拱线较远的目标，交会前必须先把漂移率降下来，离开时再重新加速。第二类处在出发段和折返段。B2、B3 是最先观测的两颗目标，R1 要从几乎不漂移的状态建立起向西的漂移；B26 和 B17 在折返之后，R1 要先把向西的漂移停下来，再建立向东的漂移。第三类是 B23 和 B27，它们与前一次观测相隔不到一天，可用于调整相位的时间很短。B19 与 B20 的经度只差 $0.03^\circ$，穿过 R1 轨道面的时刻只差 11~min，两次观测共用一组机动，为 B19 节省了一次单独的转移。

\begin{figure}[!htb]
  \centering
  \begin{subfigure}[b]{0.49\textwidth}
    \centering
    \includegraphics[width=\linewidth]{fig10a_q2_route.pdf}
    \caption{R1 相对经度随时间的变化}
  \end{subfigure}\hfill
  \begin{subfigure}[b]{0.49\textwidth}
    \centering
    \includegraphics[width=\linewidth]{fig10b_q2_dv.pdf}
    \caption{各次观测的速度增量}
  \end{subfigure}
  \caption{问题二的巡游方案}
  \label{fig:q2route}
\end{figure}

为了说明三个层次各自的作用，表~\ref{tab:q2ablation} 列出了求解过程中各阶段的结果。逐段求解得到的初始方案观测 30 颗目标，累计 210.39~m/s。访问顺序和观测时刻都不变，只做联合优化，累计速度增量降到 180.76~m/s，下降 14.1\%，这部分收益完全来自各段之间的协调。大邻域搜索把逐段求解的代价从 210.39~m/s 降到 201.74~m/s，再做联合优化，得到最终的 176.31~m/s，比初始方案下降 16.2\%。可见在本题中，联合优化的作用大于访问顺序的调整，这与第 2.3 节中“逐段独立优化不是整体最优”的判断一致。

\begin{table}[!htb]
  \centering
  \caption{问题二各求解阶段的结果}
  \label{tab:q2ablation}
  \zihao{5}
  \begin{tabular}{lccc}
    \toprule
    求解阶段 & 观测目标数 & 累计速度增量 / (m/s) & 相对初始方案 \\
    \midrule
    排序：束搜索与逐段求解 & 30 & 210.39 & — \\
    排序 + 联合优化 & 30 & 180.76 & $-14.1\%$ \\
    排序 + 改进（逐段求解） & 30 & 201.74 & $-4.1\%$ \\
    排序 + 改进 + 联合优化 & 30 & 176.31 & $-16.2\%$ \\
    \bottomrule
  \end{tabular}
\end{table}

方案的约束检验见表~\ref{tab:q2check}。优化时在距离和角度上各留 0.05~km 和 $0.05^\circ$ 的裕量，检验时各留 0.01~km 和 $0.01^\circ$，两者的观测时长都取 60.5~s，比题目要求多出 0.5~s。用解析传播在 0.05~s 网格上重放整条轨迹，并用 ATK 软件对 R1 分段接力做数值积分，两种方式都表明 30 次观测全部满足约束。全程 R1 的最低地心距为 41748.5~km，始终处在同步带附近。

\begin{table}[!htb]
  \centering
  \caption{问题二约束检验}
  \label{tab:q2check}
  \zihao{5}
  \begin{tabular}{lcc}
    \toprule
    检验项 & 要求 & 结果 \\
    \midrule
    观测目标数 & 尽可能多 & 30 / 30 \\
    累计速度增量 & $\le 300$~m/s & 176.31~m/s \\
    单次速度增量 & $\le 30$~m/s & 最大 25.71~m/s \\
    相邻脉冲间隔 & $\ge 3600$~s & 最小 3600.0~s \\
    每次观测前机动次数 & $\le 2$ & 最多 2 \\
    观测距离 & $10 \sim 50$~km & $10.07 \sim 49.95$~km \\
    最大光照角 / 侦测角 & $\le 45^\circ$ & $44.95^\circ$ / $44.95^\circ$ \\
    ATK 数值积分复核 & 全部满足 & 30 / 30 \\
    \bottomrule
  \end{tabular}
\end{table}
